% The argument behind the twenty genus-three models of data/models added on branch models-1326-210
% (PR #128): branch points located by the explicit Gross-Zagier formula, covers built exactly,
% alternatives excluded by the trace formula.  Code: HeegnerHeights.m, ExactDoubleCover.m,
% tools/heegner_model.m; the per-curve record is the PROVENANCE.md section of the same name.
%
% Numbering is local to this file.
\documentclass[11pt,reqno]{amsart}
\usepackage[margin=1.1in]{geometry}
\usepackage{amsmath,amssymb,amsthm,mathtools}
\usepackage{booktabs,array,longtable}
\usepackage[colorlinks=true,linkcolor=blue,citecolor=blue,urlcolor=blue]{hyperref}
\theoremstyle{plain}
\newtheorem{theorem}{Theorem}
\newtheorem{proposition}[theorem]{Proposition}
\newtheorem{lemma}[theorem]{Lemma}
\theoremstyle{definition}
\newtheorem{remark}[theorem]{Remark}
\DeclareMathOperator{\Jac}{Jac}
\DeclareMathOperator{\Gal}{Gal}
\DeclareMathOperator{\Div}{div}
\DeclareMathOperator{\Area}{Area}
\DeclareMathOperator{\ord}{ord}
\newcommand{\Q}{\mathbb{Q}}
\newcommand{\Z}{\mathbb{Z}}
\newcommand{\F}{\mathbb{F}}
\newcommand{\Qbar}{\overline{\mathbb{Q}}}
\newcommand{\hh}{\hat h}
\begin{document}
\title{Twenty genus-three Shimura curve quotients\\ from the heights of their branch points}
\author{}
\date{\today}
\maketitle

\section{The statement}

Let $X=X_0^D(N)/W$ be an Atkin--Lehner quotient of genus three that has an elliptic quotient
$E=X/\langle W,w\rangle$ by one further Atkin--Lehner involution. Then $X\to E$ is a double cover
branched at four points, and those points are CM points whose discriminants are known. The explicit
Gross--Zagier formula of Cai, Shu and Tian gives the canonical heights of their images on $E$; the
points of $E$ of those heights are finitely many; every double cover of $E$ over $\Q$ branched
along each candidate divisor is constructed exactly, and all but one are excluded by their point
counts against the trace formula, which needs no model. This note writes out that argument and
records the twenty curves it identifies. They are the genus-three quotients of the classification
that lie beyond the reach of the Borcherds-product constructions, each being a double cover of a
genus-one curve rather than of $\mathbb P^1$.

\begin{theorem}\label{thm:main}
Each quotient $X=X_0^D(N)/W$ of Table~\ref{tab:models} is $\Q$-isomorphic to the curve listed
beside it, with the following qualifications. For $X_0^6(385)/W$ and $X_0^{714}(5)/W$ the
identification is conditional on the completeness of the enumeration over the isogeny class of $E$,
which Magma's Mordell--Weil computation does not certify on every curve of the class; for
$X_0^{15}(38)/W$ and $X_0^6(323)/W$ it is conditional on a finite search window for one branch place
whose position the formula does not predict. In every case the listed curve is a double cover of $E$
branched at the CM points of the stated discriminants, is hyperelliptic of genus three, and has the
zeta function of $X$ at every good prime below $200$ and over $\F_{p^2}$ at every good prime below
$40$.
\end{theorem}

Three of the twenty have no model $y^2=f(x)$ over $\Q$: their canonical image is the conic
$X_1^2+2X_2^2+2bX_3^2=0$ ramified at $2$ and at infinity, which has no rational point to project
from, and they are given as double covers of that conic. The thirty-fourth open quotient,
$X_0^{141}(2)/\{1,6,94,141\}$, is decided by the same computation in the other direction: the one
configuration consistent with the trace formula there is not hyperelliptic.

\begingroup
\footnotesize
\setlength{\tabcolsep}{5pt}
\begin{longtable}{@{}p{4.6cm} p{1.6cm} p{8.9cm}@{}}
\caption{The twenty quotients: $W$ by generators, $E$ by Cremona label, the branch places by the
CM discriminant of their points (Ogg), with the degree of the place when it is not one and its
field; below each, the model. $\dagger$: the quaternionic form is old at $2$, so $E$ has conductor
$DN/2$. For the conic models $X_1,X_2,X_3$ have weight $1$ and $y$ weight $2$; $s,z$ stand for
$X_1,X_2$.}\label{tab:models}\\
\toprule
quotient & $E$ & branch places \\
\midrule
\endfirsthead
\toprule
quotient & $E$ & branch places \\
\midrule
\endhead
$X_0^{1326}(1)/\langle w_{17},w_{78}\rangle$ & 1326b2 & $-24$, $-408$: both of degree 2 over $\Q(\sqrt2)$\\
\multicolumn{3}{@{}l@{}}{\quad$\scriptstyle y^2=272x^8-72x^6-55x^4-18x^2+17$}\\[2pt]
$X_0^{210}(13)/\langle w_2,w_{21},w_{65}\rangle$ & 2730u1 & $-420$, $-1092$: degree 2 over $\Q(i)$ (also from 2730s1: $-312$, $-120$ over $\Q(\sqrt2)$)\\
\multicolumn{3}{@{}l@{}}{\quad$\scriptstyle y^2=4x^8-20x^7+97x^6-20x^5-138x^4+20x^3+97x^2+20x+4$}\\[2pt]
$X_0^{39}(11)/\langle w_{11},w_{39}\rangle$ & 429b1 & $-52$ over $\Q(\sqrt{13})$, $-132$ over $\Q(\sqrt{-11})$\\
\multicolumn{3}{@{}l@{}}{\quad$\scriptstyle y^2=-891x^8+1044x^6+214x^4-124x^2+13$}\\[2pt]
$X_0^{55}(7)/\langle w_{35},w_{55}\rangle$ & 385a1 & $-20$ over $\Q(\sqrt5)$, $-1540$ over $\Q(\sqrt{-35})$\\
\multicolumn{3}{@{}l@{}}{\quad$\scriptstyle y^2=-35x^8-2028x^6+198x^4+68x^2+5$}\\[2pt]
$X_0^{6}(301)/\langle w_3,w_{14},w_{43}\rangle$ & 1806l1 & $-1032$ over $\Q(\sqrt{129})$, $-1204$ over $\Q(\sqrt{-7})$\\
\multicolumn{3}{@{}l@{}}{\quad$\scriptstyle y^2=129x^8-2276x^6+3334x^4+1180x^2-63$}\\[2pt]
$X_0^{6}(455)/\langle w_6,w_7,w_{10},w_{13}\rangle$ & 2730y1 & $-1560$ over $\Q(\sqrt{10})$, $-10920$ over $\Q(\sqrt{-26})$\\
\multicolumn{3}{@{}l@{}}{\quad$\scriptstyle y^2=-416x^8-515x^6+582x^4-251x^2+40$}\\[2pt]
$X_0^{6}(187)/\langle w_3,w_{11},w_{17}\rangle$ & 1122i1 & $-264$, $-4488$: both of degree 2 over $\Q(\sqrt{-2})$\\
\multicolumn{3}{@{}l@{}}{\quad$\scriptstyle X_1^2+2X_2^2+22X_3^2=0,\ \ y^2=4449s^4+84056s^3z+411472s^2z^2+195328sz^3+797696z^4$}\\[2pt]
$X_0^{22}(51)/\langle w_6,w_{17},w_{22}\rangle$ & 1122f1 & $-264$, $-4488$: both of degree 2 over $\Q(\sqrt{-2})$\\
\multicolumn{3}{@{}l@{}}{\quad$\scriptstyle X_1^2+2X_2^2+59X_3^2=0,\ \ y^2=9242176s^4-2871136s^3z+4331508s^2z^2+124292sz^3+4787449z^4$}\\[2pt]
$X_0^{6}(247)/\langle w_2,w_{19},w_{39}\rangle$ & 1482i1 & $-3$, $-52$ rational; $-456$ of degree 2 over $\Q(\sqrt{57})$\\
\multicolumn{3}{@{}l@{}}{\quad$\scriptstyle y^2=81x^8-4788x^6+1654x^4+2540x^2+513$}\\[2pt]
$X_0^{6}(253)/\langle w_2,w_{23},w_{33}\rangle$ & 1518p1 & $-88$, $-1012$ rational; $-552$ of degree 2 over $\Q(\sqrt{-23})$\\
\multicolumn{3}{@{}l@{}}{\quad$\scriptstyle y^2=9x^8+1388x^6+3526x^4-3060x^2-1863$}\\[2pt]
$X_0^{6}(185)/\langle w_2,w_{15},w_{111}\rangle$ & 1110m1 & $-40$, $-555$ rational; $-4440$ of degree 2 over $\Q(\sqrt{-111})$\\
\multicolumn{3}{@{}l@{}}{\quad$\scriptstyle y^2=2025x^8+4716x^6-122x^4-5620x^2-999$}\\[2pt]
$X_0^{6}(385)/\langle w_5,w_6,w_{11},w_{14}\rangle$ & 2310q1 & $-40$, $-1540$ rational; $-9240$ of degree 2\\
\multicolumn{3}{@{}l@{}}{\quad$\scriptstyle y^2=9x^8-580x^6-9482x^4+60092x^2+3465$} \hfill\textit{candidate}\\[2pt]
$X_0^{714}(5)/\langle w_3,w_{10},w_{14},w_{17}\rangle$ & 3570x2 & $-340$, $-840$ rational; $-14280$ of degree 2\\
\multicolumn{3}{@{}l@{}}{\quad$\scriptstyle y^2=9720x^8+2061x^6+598x^4-155x^2+16$} \hfill\textit{candidate}\\[2pt]
$X_0^{15}(38)/\langle w_2,w_{15},w_{19}\rangle$ & 570e4 & $-12$, $-40$, $-228$, $-760$, all rational\\
\multicolumn{3}{@{}l@{}}{\quad$\scriptstyle y^2=81x^8+180x^6+2062x^4-1660x^2+361$} \hfill\textit{window at $-12$}\\[2pt]
$X_0^{141}(2)/\langle w_3,w_{94}\rangle$ & 141a1$^{\dagger}$ & $-4$, $-24$ rational; $-564$ of degree 2 over $\Q(\sqrt{-47})$\\
\multicolumn{3}{@{}l@{}}{\quad$\scriptstyle y^2=81x^8+252x^6-26x^4-260x^2-47$}\\[2pt]
$X_0^{21}(34)/\langle w_3,w_{14},w_{17}\rangle$ & 357d1$^{\dagger}$ & $-4$, $-84$, $-408$, $-1428$, all rational\\
\multicolumn{3}{@{}l@{}}{\quad$\scriptstyle y^2=81x^8-612x^6+262x^4+220x^2+49$}\\[2pt]
$X_0^{57}(10)/\langle w_5,w_6,w_{19}\rangle$ & 285a1$^{\dagger}$ & $-4$, $-760$ rational; $-1140$ of degree 2\\
\multicolumn{3}{@{}l@{}}{\quad$\scriptstyle y^2=729x^8+500x^6-2x^4+68x^2-15$}\\[2pt]
$X_0^{15}(46)/\langle w_5,w_6,w_{23}\rangle$ & 345f1$^{\dagger}$ & $-15$, $-40$, $-60$, $-1380$, all rational\\
\multicolumn{3}{@{}l@{}}{\quad$\scriptstyle y^2=81x^8-900x^6+358x^4+380x^2+81$}\\[2pt]
$X_0^{690}(7)/\langle w_6,w_{10},w_{14},w_{23}\rangle$ & 4830bj1 & $-3$, $-115$, $-280$, $-420$, all rational\\
\multicolumn{3}{@{}l@{}}{\quad$\scriptstyle y^2=9x^8-130x^6-767x^4-520x^2+144$}\\[2pt]
$X_0^{6}(323)/\langle w_6,w_{19},w_{34}\rangle$ & 1938h1 & $-7752$: one place of degree 4\\
\multicolumn{3}{@{}l@{}}{\quad$\scriptstyle X_1^2+2X_2^2+22X_3^2=0,\ \ y^2=46521s^4+32184s^3z+101968s^2z^2+91392sz^3+41984z^4$} \hfill\textit{window}\\[2pt]
\bottomrule
\end{longtable}
\endgroup

\section{The branch locus}\label{sec:branch}

Write $W'=\langle W,w\rangle$, so that $E=X_0^D(N)/W'$ has genus one, and $\pi\colon X\to E$ is
the quotient by the involution that $W'/W$ induces on $X$. A point of $E$ is a branch point of
$\pi$ exactly when its $W'$-orbit on $X_0^D(N)$ is a single $W$-orbit, that is, when the orbit
consists of fixed points of some $w_m$ with $m\in W'\setminus W$. By Ogg \cite{ogg} the fixed
points of $w_m$ on $X_0^D(N)$ are the CM points by the quadratic orders $R$ of discriminant $-4m$
and, for $m\equiv3\pmod 4$, $-m$ (and for $m=2$ also $-4$, $-8$), and $R$ contributes
$h(R)\prod_{p\mid DN,\,p\nmid m}\bigl(1-\{R/p\}\bigr)$ of them, $\{R/p\}$ the Eichler symbol. The
count divided by $|W|$ is the number of points of $X$ over the branch points of that discriminant,
and Riemann--Hurwitz for $\pi$, $2\cdot3-2=2\cdot0+\deg B$, says the branch divisor $B$ has degree
four. The fields of definition of the points of $X$ over a branch point of discriminant $d$ come
from Shimura reciprocity: the CM points of discriminant $d$ are defined over the ring class field,
and the Galois action on them is that of the class group twisted by the Atkin--Lehner action of
the primes dividing $DN$ (\texttt{BranchPointData}). In the twenty cases the places of $B$ have
degree one or two, except at $X_0^6(323)/W$, where $B$ is a single place of degree four. The third
and fourth columns of Table~\ref{tab:models} are this computation.

\section{The heights of the branch points}\label{sec:heights}

Let $f\colon X_0^D(N)\to E$ be the composite $X_0^D(N)\to X\to E$, of degree $2|W|$, and let $P$ be
a CM point of fundamental discriminant $d$ on $X_0^D(N)$, with $h=h(d)$ and $u=[\mathcal
O_d^\times:\Z^\times]$. For a genus character $\chi=(d_1,d_2)$ of the class group of $d$, that is a
factorisation $d=d_1d_2$ into fundamental discriminants, put $L(s,E,\chi)=L(E^{d_1},s)L(E^{d_2},s)$
and let $P^0_\chi=\sum_{t}\chi(t)\,f(P)^t$, the sum over the class group acting through
Shimura reciprocity. Theorem 1.5 of Cai, Shu and Tian \cite{cst}, in its second special case
($F=\Q$, trivial central character, a character of conductor $1$), states
\[
  L'(1,E,\chi)\;=\;2^{-\#\Sigma_D}\,\frac{8\pi^2(\varphi,\varphi)_{\Gamma_0(N_E)}}{u^2\sqrt{|d|}}\,
  \frac{\langle P^0_\chi,P^0_{\bar\chi}\rangle_K}{\deg f},
\]
where $\varphi$ is the normalised newform of $E$, $\Sigma_D$ is the set of primes dividing both
$N_E$ and $d$, and the pairing is the N\'eron--Tate pairing over $K=\Q(\sqrt d)$, which is twice
the absolute height. The theorem's test-vector condition, in our situation, is the following:
at a prime $v\in\Sigma_D$ the involution induced by a uniformiser of $K_v$ is the Atkin--Lehner
involution $w_v$, and the formula holds for the characters $\chi$ whose value at the ramified
prime above $v$ equals the eigenvalue $\varepsilon_v$ of $w_v$ on the quaternionic form, which is
$a_v(E)$ for $v\mid D$ and $-a_v(E)$ for $v\mid N$ (the Jacquet--Langlands correspondence changes the
sign at the primes of $D$); for the other characters the component $P^0_\chi$ is torsion. We call
the former characters \emph{admissible}. If $\pi\colon X_0(N_E)\to E_{\mathrm{opt}}$ is the optimal
parametrisation then $\pi^*\omega=2\pi i\,\varphi(z)\,dz$ up to the Manin constant, which is $1$ for
squarefree conductor \cite{cesnavicius}, so $8\pi^2(\varphi,\varphi)=2\deg\pi\cdot\Area(E_{\mathrm
{opt}})$, and the formula becomes
\begin{equation}\label{eq:height}
  \hh(P^0_\chi)\;=\;\frac12\cdot\frac{u^2\sqrt{|d|}\,\deg f\;2^{\#\Sigma_D}\,L'(1,E,\chi)}
  {2\deg\pi\cdot\Area(E_{\mathrm{opt}})}
\end{equation}
for the absolute canonical height (\texttt{HeegnerHeightComponents}). Since $f(P)=\frac1h\sum_\chi
P^0_\chi$ and the $P^0_\chi$ for distinct characters are orthogonal,
$\hh(f(P))=h^{-2}\sum_{\chi\ \mathrm{adm.}}\hh(P^0_\chi)$ when the class group is $2$-elementary.
What is used, though, is finer than the height of $f(P)$. A branch place of degree two with field
$F=\Q(\sqrt{d_F})$ is a pair $\{f(P),\overline{f(P)}\}$ of conjugate points, and the only
characters with a nonzero component on a point defined over $F$ are the trivial one and the genus
character $\chi_F$ attached to $F$; so
\[
  f(P)-\overline{f(P)}=\tfrac{2}{h}\,P^0_{\chi_F}\in E(F)^-,\qquad
  f(P)+\overline{f(P)}=\tfrac{2}{h}\,P^0_{1}+(\text{torsion})\in E(\Q),
\]
and \eqref{eq:height} gives the height of each side exactly. The twists $E^{d_1}$, $E^{d_2}$ that
occur have analytic rank at most one, and the L-values are computed numerically; a nonzero value
is recognised against a threshold, so a value wrongly taken for zero can only remove candidates
and never admit one.

\subsection*{The rank-one shift.}
When $E$ has rank one over $\Q$ the trivial-character component is not given by
\eqref{eq:height} alone: on a modular curve the Heegner divisor must be made of degree zero, and
the image of the Hodge class of $X_0^D(N)$ enters. Concretely, the sum of the points of $X$ over a
branch place of degree $e$ is
\[
  e\,\xi_E\;\pm\;b\,G\;+\;(\text{torsion}),
\]
with $\xi_E\in E(\Q)\otimes\Q$ the image of the Hodge class, the same for every place, $G$ a
generator of $E(\Q)$ modulo torsion, and $b$ a rational number whose absolute value
\eqref{eq:height} gives (as $(2/h)\sqrt{\hh(P^0_1)/\hh(G)}$; zero when the trivial character is
inadmissible). This is what the comparison with the branch points of the models already known
showed: $82$ of $92$ branch points matched \eqref{eq:height} to six digits, and the ten that did
not were all on curves of rank one, shifted by a common amount. Translating every branch point by
one rational point of $E$ gives an isomorphic cover, so only the class of $\xi_E$ modulo $\Z G$
(and torsion) matters, and its denominator divides twice the denominators of the $b$'s; the
enumeration over the signs of the $b$'s and over that finite set of shifts is therefore finite and
complete (\texttt{ModelByHeegnerHeights}). At $X_0^6(253)/W$ the shift has denominator $3$, and an
earlier version of the search that assumed it integral found nothing.

\subsection*{Oldforms.}
At five of the twenty the quaternionic form is old at $p=2$: $E$ has conductor $DN/2$ and the
routine that reads an elliptic quotient off the Atkin--Lehner signs does not find it, so $E$ was
identified from the trace formula by matching the point counts of the genus-one quotient against
the elliptic curves of conductor dividing $DN$. The map to which the formula applies is then
$g=f\pi_1+f\pi_2$ with $\pi_1,\pi_2$ the two degeneracy maps, $\pi_2=\pi_1w_p$, and $f$ the
parametrisation from the curve on which the form is new; its degree is $2(p+1+a_p)\deg f$, each
$f\pi_i$ having degree $(p+1)\deg f$ and each cross term $fT_pf^\vee$ degree $a_p\deg f$. Where $p$
ramifies in the field of the CM point, $w_p$ acts on it as the Frobenius of the prime above $p$,
so $\pi_2(P)$ is that conjugate of $\pi_1(P)$ and the $\chi$-component of $g(P)$ is $(1+\chi(p))$
times that of $f(\pi_1P)$. With $\deg g$ in place of $\deg f$ in \eqref{eq:height}, every
component is multiplied by
\[
  \frac{(1+\chi(p))^2}{2(p+1+a_p)},
\]
which is $2/(p+1+a_p)$ for $\chi(p)=1$ and $0$ for $\chi(p)=-1$: $2$ at 141a1, 357d1 and 345f1
($a_2=-2$) and $1$ at 285a1 ($a_2=-1$). The vanishing at $\chi(p)=-1$ is what accounts for the
rational branch places whose nontrivial characters would otherwise have had nonzero components.
This derivation is Sachi Hashimoto's. It is for $p$ ramified in the field of the branch point; the
same expression is applied at a place where $p$ is not ramified, which is not derived, and there
the height only directs the search while the counts certify the answer.

\section{From heights to points}\label{sec:points}

Fix $E$ in its isogeny class. For a branch place of degree two with field $F$, the minus part
$f(P)-\overline{f(P)}$ is a point of $E(F)^-$, which is the Mordell--Weil group of the quadratic
twist $E^{d_F}$; in the cases at hand that group has rank one, and Magma proves its rank and
saturation (with \texttt{Effort := 2} at $X_0^{714}(5)/W$). The height of the minus part is known
exactly, so it is $\pm k\,G_F$ for a determined integer $k$ --- and that $k$ is an integer is a
test: on the other curves of the isogeny class the predicted height is not the height of a lattice
point, which is how $E$ is identified within its class (at $X_0^{39}(11)/W$ it excludes five of the
six curves of 429b; at $X_0^6(455)/W$ two curves pass this necessary test and are excluded only by
the counts). The plus part is $2\xi_E\pm bG$ as above. A place of degree one is a point of $E(\Q)$
equal to $\xi_E\pm bG$ up to torsion; when every place is rational there is no minus part and no
lattice test, and the integrality of the rational parts takes its place: at $X_0^{15}(38)/W$ it
excludes three of the four curves of 570e. Two places are outside the formula as implemented: a
discriminant that is not fundamental ($-12$ at $X_0^{15}(38)/W$, the order of conductor $2$ in
$\Q(\sqrt{-3})$), whose multiple of $G$ is searched over $|n|\le12$ instead; and the place of
degree four at $X_0^6(323)/W$, whose pinning characters have order four. There the trivial
character gives the sum of the four points as $4\xi_E\pm2G$, the character of the quadratic
subfield $\Q(\sqrt{-2})$ of the quartic field has a vanishing component because the twist of 1938h1
by $-8$ has analytic rank two, and the one unknown left is the minus part for the quartic field
over $\Q(\sqrt{-2})$, a multiple of the generator of a rank-one Mordell--Weil group (proved and
saturated, height $7.1168638516$); multiples $|k|\le4$ were searched and $k=\pm1$ is the only
survivor. These two are the window-conditional cases of Theorem~\ref{thm:main}.

The outcome of this section is, for each $E$ in the isogeny class, a finite list of candidate
branch divisors $B$, each a reduced $\Q$-rational divisor of degree four on $E$, and the guarantee
that the true branch divisor is among them (within the stated windows).

\section{The covers, and the certification}\label{sec:covers}

A double cover of $E$ over $\Q$ branched along a reduced $\Q$-rational divisor $B$ of degree four
has function field $\Q(E)(\sqrt\gamma)$ with $\Div\gamma=B-2\Delta$, $\Delta$ of degree two. Writing
$\Delta\sim(Q)+(O)$, the condition $2\Delta\sim B$ says $2Q=P_B$, where $P_B$ is the point with
$B\sim(P_B)+3(O)$; so the covers over $\Q$ are indexed by the rational halves $Q$ of $P_B$, a torsor
under $E(\Q)[2]$ with at most four elements, and each is determined up to a quadratic twist along
the deck involution, $\gamma\mapsto\kappa\gamma$ with $\kappa\in\Q^\times/\Q^{\times2}$
(\texttt{DoubleCoversWithBranchDivisor}). Whether such a cover is hyperelliptic is read off its
function field: in genus three the pairwise products of a basis of the holomorphic differentials
span a space of dimension $5$ if the curve is hyperelliptic, its canonical image being a conic, and
$6$ otherwise. When the conic has a rational point, projecting from it gives a function $t$ of
degree two and a model $y^2=G(t)$; when it has none, the curve has no such model over $\Q$ and is
written as $aX_1^2+bX_2^2+cX_3^2=0$, $y^2=G(X_1,X_2,X_3)$ (\texttt{HyperellipticModelOverConic}),
whose points over $\F_q$ are counted by parametrising the conic, which is isotropic over every
finite field (\texttt{PointCountOverConic}).

\begin{proposition}\label{prop:certify}
Let $X\to E$ be as above, let $\mathcal B$ be a finite set of reduced $\Q$-rational divisors of
degree four on $E$ containing the branch divisor of $X$, and for each $B\in\mathcal B$ let
$\mathcal C_B$ be the covers over $\Q$ branched along $B$, up to twist. Suppose that for every
$(B,C)$ with $C\in\mathcal C_B$ except one, there is a good prime $p$ at which $\#X(\F_p)$, given by
the trace formula, equals neither $\#C(\F_p)$ nor the count of the twist of $C$ by a non-square
modulo $p$; and that for the remaining $(B_0,C_0)$ exactly one squarefree $\kappa$ supported on
$\{-1,2\}$, the primes of bad reduction of $X$, and the primes where $C_0$ can have bad reduction
(those dividing the discriminant of $E$ or a coefficient of $\gamma$, and those modulo which two
points of $B_0$ coincide) has $\#C_0^{\kappa}(\F_p)=\#X(\F_p)$ at every good prime used. Then
$X\cong C_0^{\kappa}$ over $\Q$.
\end{proposition}

\begin{proof}
$X$ is a double cover of $E$ over $\Q$ branched along its branch divisor $B\in\mathcal B$, hence
$X\cong C^{\kappa'}$ for some $C\in\mathcal C_B$ and some $\kappa'$. At a good odd prime $p$ the
reductions of $C$ and of its twist by a non-square have the two possible counts of any twist
$C^{\kappa'}$, so a pair $(B,C)$ matching neither at some $p$ is not $X$; this leaves $(B_0,C_0)$.
A twist $\kappa'$ ramified at an odd prime $p$ where $C_0$ has good reduction gives $C_0^{\kappa'}$
bad reduction at $p$, and $X$ has good reduction at every prime not dividing $DN$ (Morita, and the
quotient by a group of order prime to $p$), so $\kappa'$ is supported on the listed primes; the
counts at the good primes then determine the Kronecker symbols of $\kappa'$ at each of them, and by
hypothesis only $\kappa$ fits. (\texttt{MatchCoverToCounts}; a cover whose counts match one of its
two twists at every prime separately but whose signs are the symbols of no admissible $\kappa$ is
excluded by the same rule.)
\end{proof}

Theorem~\ref{thm:main} follows: Sections~\ref{sec:branch}--\ref{sec:points} produce $\mathcal B$ with
the guarantee that it contains the branch divisor --- unconditionally where every place is
predicted by \eqref{eq:height} and the Mordell--Weil groups are proved, within the window at
$X_0^{15}(38)/W$ and $X_0^6(323)/W$, and conditionally on the isogeny class having been exhausted at
$X_0^6(385)/W$ and $X_0^{714}(5)/W$ --- and \texttt{tools/heegner\_model.m} runs the proposition's
test over all of it, returning one curve up to $\Q$-isomorphism in each case. Where two
configurations survive (twists $-14$ and $721$ at $X_0^6(187)/W$, $-43$ and $-2171$ at
$X_0^{714}(5)/W$, and the two presentations at $X_0^6(323)/W$) the isomorphism between the resulting
curves is exhibited, not inferred from the counts. That the selection by counts is part of the
argument and not a formality is shown by $X_0^6(23)/w_{23}$, whose $E$ has $E(\Q)[2]$ of order two:
there are two covers over $\Q$ with the same $E$ and the same four branch points, one hyperelliptic
and one not, and the trace formula is consistent with the second only.

\section{What was predicted, and what rests on what}\label{sec:checks}

The following were not inputs to the construction and came out as predicted for every model: the
genus is three; the primes of bad reduction are exactly the primes of $DN$ (for the conic models,
of the two equations' discriminants, apart from a $59$ at $X_0^{22}(51)/W$ that is an artefact of
the coordinates, the conic having points over every $\Q_p$ with $p$ odd); the point counts agree
with the trace formula at every good prime below $200$ and over $\F_{p^2}$ at every good prime
below $40$, which in genus three determines the zeta function; the quotient by the residual
involution is $E$; the automorphism group has order $4$, or $8$ at $X_0^{1326}(1)/W$ and
$X_0^{210}(13)/W$ where $W$ has index four in the Atkin--Lehner group, with the three elliptic
quotients the signs predict (1326b2, 1326c1, 1326d1; 2730u1, 2730s1, 210a2); and the fields of
definition of the fixed points of the involutions are the fields of the branch places, rational
exactly at the places of degree one. The models of $X_0^{39}(11)/W$ and $X_0^{55}(7)/W$ are
$\Q$-isomorphic to the curves fitted numerically from the Shimizu-lifted differentials, a route
sharing no step with this one; at $X_0^6(185)/W$, where that fit could not recognise two of its
invariants, the exact model agrees with the numerical values to twelve digits.

What the theorem rests on: Theorem 1.5 of \cite{cst} with the admissibility rule, which is its own
test-vector condition; the Manin constant $1$ for squarefree conductor \cite{cesnavicius}, recorded
as $1$ in the LMFDB for every curve used; Magma's proofs of the Mordell--Weil ranks and saturations
(the two candidates are exactly where that proof is missing on part of the isogeny class); the
oldform multiplier, derived where $p$ ramifies in the field of the branch point; the enumeration of
Section~\ref{sec:points}, complete except within the two windows; the exact double-cover
construction of Section~\ref{sec:covers}; and the trace formula for the counts of $X$. The
L-values are the only numerical quantities, and they direct the search rather than certify it.

\begin{remark}
The candidates can be settled without new Mordell--Weil computations. If $\phi\colon E\to E'$ is an
isogeny of degree $m$ inside the class, the generator of the twist of $E'$ is $(c/m)\phi(G)$ for an
integer $1\le c\le m$, so its height is $c^2/m$ times the known one, and the height comparison on
$E'$ can be run over those finitely many values.
\end{remark}

\section{The computations}\label{sec:computations}

\begin{itemize}
\item \texttt{HeegnerHeights.m}: \texttt{QuaternionicAtkinLehnerSign}, \texttt{EllipticQuotientsBySigns},
  \texttt{BranchPointData} (Section~\ref{sec:branch}), \texttt{GenusCharacters},
  \texttt{HeegnerHeightComponents} and \texttt{PredictedBranchPointHeight} (Section~\ref{sec:heights}),
  \texttt{ModelByHeegnerHeights} (Section~\ref{sec:points} and the search), \texttt{SmallHyperellipticModel},
  \texttt{KleinFourEvenModel}.
\item \texttt{ExactDoubleCover.m}: \texttt{DoubleCoversWithBranchDivisor}, \texttt{TwistedCover},
  \texttt{HyperellipticModelOfCover}, \texttt{HyperellipticModelOverConic}, \texttt{PointCountOverConic},
  \texttt{CoverPointCounts}, \texttt{MatchCoverToCounts} (Section~\ref{sec:covers}).
\item \texttt{tools/heegner\_model.m}: the search for a given $(D,N,W,w)$ over the whole isogeny
  class; the command for each of the twenty is in the \texttt{PROVENANCE.md} section ``Twenty
  quotients from the heights of their branch points'', which also records each run's survivors.
\item \texttt{tests/HeegnerHeights.m}: the formula against the branch points of the stored models
  of $X_0^{35}(1)$, $X_0^{10}(11)/\{1,5\}$, $X_0^{390}(1)/\{1,3,10,30\}$, and the route end to end on
  the last two; \texttt{tests/ExactDoubleCover.m}: the covers of $X_0^{35}(1)$, $X_0^{51}(1)$,
  $X_0^6(23)/w_{23}$ and a configuration over a pointless conic; \texttt{tests/ModelChecks.m}: every
  stored model against the trace formula; \texttt{tests/CRVStructure.m}: the conic entries as
  irreducible curves of genus three.
\end{itemize}

\begin{thebibliography}{9}
\bibitem{cst} L.~Cai, J.~Shu and Y.~Tian, \emph{Explicit Gross--Zagier and Waldspurger formulae},
Algebra Number Theory \textbf{8} (2014), 2523--2572.
\bibitem{cesnavicius} K.~\v{C}esnavi\v{c}ius, \emph{The Manin constant in the semistable case},
Compos. Math. \textbf{154} (2018), 1889--1920.
\bibitem{gz} B.~H. Gross and D.~B. Zagier, \emph{Heegner points and derivatives of $L$-series},
Invent. Math. \textbf{84} (1986), 225--320.
\bibitem{ogg} A.~P. Ogg, \emph{Real points on Shimura curves}, in: Arithmetic and geometry, Vol.~I,
Progr. Math. \textbf{35}, Birkh\"auser, 1983, 277--307.
\bibitem{yzz} X.~Yuan, S.-W. Zhang and W.~Zhang, \emph{The Gross--Zagier formula on Shimura curves},
Ann. of Math. Stud. \textbf{184}, Princeton, 2013.
\end{thebibliography}
\end{document}
